\section{Exercises}
\label{sec:fol-cardinality-normalforms-06}

\begin{exercise}
\label{ex:fol-cardinality-normalforms-01}
Prove the Tarski--Vaught test in full detail, handling both quantifiers
directly in the structural induction and justifying every appeal to the
induction hypothesis.
\end{exercise}

\begin{exercise}
\label{ex:fol-cardinality-normalforms-02}
Let \(\Meta\) be a structure in a countable signature and let
\(A\subseteq\abs{\Meta}\) be countable.  Write out the first three stages
\(A_0,A_1,A_2\) of the downward L\"owenheim--Skolem construction and prove
that the union \(\bigcup_{n\in\N}A_n\) of all stages is closed under every
signature function.
\end{exercise}

\begin{exercise}
\label{ex:fol-cardinality-normalforms-03}
Use compactness to prove that if a theory has arbitrarily large finite models,
then it has an infinite model.  State why this conclusion alone does not say
that any finite model embeds into the infinite model.
\end{exercise}

\begin{exercise}
\label{ex:fol-cardinality-normalforms-04}
Prove that no first-order sentence in the empty signature is true exactly in
the finite nonempty structures.
\end{exercise}

\begin{exercise}
\label{ex:fol-cardinality-normalforms-05}
Convert
\[
 \neg\forall x\bigl(P(x)\to\exists y(Q(y)\land R(x,y))\bigr)
\]
to rectified prenex normal form.  Justify each equivalence and side condition.
\end{exercise}

\begin{exercise}
\label{ex:fol-cardinality-normalforms-06}
Skolemize
\[
 \exists u\forall x\exists y\forall z\exists w\,
 R(u,x,y,z,w).
\]
Give the arity and argument list of every new symbol.
\end{exercise}

\begin{exercise}
\label{ex:fol-cardinality-normalforms-07}
Give two structures showing why a sentence and its Skolemization cannot be
compared for logical equivalence before an expanded signature and a reduct map
have been specified.
\end{exercise}

\begin{exercise}
\label{ex:fol-cardinality-normalforms-08}
Find a most general unifier, or prove none exists, for each pair:
\[
 P(f(x),x),P(y,g(y));
 \qquad
 P(f(x),x),P(f(y),y);
 \qquad
 P(x,x),P(y,f(y)).
\]
The first and final pairs require the occurs check.
\end{exercise}

\begin{exercise}
\label{ex:fol-cardinality-normalforms-09}
Clausify and refute by first-order resolution:
\[
 \forall x(P(x)\to Q(x)),
 \qquad
 \exists xP(x),
 \qquad
 \forall x\neg Q(x).
\]
Record the Skolem constant, every substitution, and the empty clause.
\end{exercise}

\begin{exercise}
\label{ex:fol-cardinality-normalforms-10}
Explain exactly why failure to find a resolution refutation within
a fixed time bound does not establish satisfiability.
\end{exercise}
